\chapter{構造地質}
\section{何を復元する学問か}
岩石がいつ、どこで、どのように変形したかを、現在の露頭から復元する分野である。
層理や鉱物の配列は変形前後を結ぶ標識となる。
一枚の露頭は三次元構造の一断面であり、姿勢測定・層序・微小組織・地下探査を合わせて
幾何、運動、形成条件、時間順序を推定する。
地質学的な形態名は、物理学の応力やひずみそのものではない。
形を認識した後に、どの力学モデルでその形を説明できるかを問う\cite{str_orientation}。

\section{面と線の姿勢：走向・傾斜・沈下角}
\label{theory:structure:orientation}
地層の境界や断層面は面として、条線や褶曲のヒンジは線として記載する。
面と水平面の交線の方位が走向である。最大の下降を示す方向は走向に直交し、
その鉛直断面内で水平から測る角が真の傾斜である。
走向に平行な鉛直断面では、平面との交線は水平となり、真の傾斜を測れない。
面の法線を \(\boldsymbol n\)、鉛直単位ベクトルを \(\boldsymbol e_z\) とすると、
走向の軸は \(\boldsymbol e_z\times\boldsymbol n\) で表せる。
水平面ではこの外積が零となり走向は一意でない。鉛直面では下降側の選び方にも注意する。

線の姿勢は、下向きの水平投影方位と水平からの沈下角（プランジ）で表す。
線の単位方向 \(\boldsymbol l\) では、沈下角は
\(\arcsin(|\boldsymbol l\cdot\boldsymbol e_z|)\) である。
絶対値は軸としての傾きのみを表し、変位センスの正負は別に定める。
クリノメータは水平を作る部分と鉛直面内の角を測る部分を使い分ける。
磁針を使う方位には偏角補正・近傍磁性体・計器の校正も必要である。
原典の「線構造の走向」は線の水平投影方位の呼称として保持し、面の走向とは区別する
\cite{str_orientation}。

\section{褶曲の幾何と年代}
\label{theory:structure:folds}
曲がった地層のうち曲率の大きい部分がヒンジ、比較的曲率の小さい両側が翼である。
一つの層上のヒンジを結ぶヒンジ線と、複数層のヒンジを含む軸面を分ける。
上に凸の形がアンチフォーム、下に凸の形がシンフォームである。
背斜は中心に古い層、向斜は中心に新しい層を持つという層序の概念である。
地層の逆転や再褶曲があると、上に凸という観察だけから背斜とは決められない
\cite{str_fold_geometry}。
カルミネーションは軸が延長方向に高まる部分であり、沈下角の変化と関係する。
単一断面での山形と、軸方向の高まりは別の情報である\cite{str_culmination}。

\begin{center}
\setlength{\unitlength}{1pt}
\begin{picture}(240,82)
\qbezier(20,15)(120,105)(220,15)
\qbezier(20,5)(120,95)(220,5)
\qbezier(20,-5)(120,85)(220,-5)
\put(112,63){\makebox(0,0){ヒンジ}}
\put(40,45){\makebox(0,0){翼}}
\put(200,45){\makebox(0,0){翼}}
\put(120,-5){\line(0,1){67}}
\end{picture}
\par\small 独自模式断面：上に凸の形態である。地層の年代が未指定なので背斜とは確定しない。
\end{center}

\section{座屈と曲げ：力の向きから読む}
\label{theory:structure:buckling}
層に沿う短縮により、初期の小さな曲がりが増幅するのが座屈褶曲である。
一方、層に大きな角度で荷重が作用して曲げる形式もある。
粘性の異なる層と周囲の物質は変形を分担し、層の厚さ・物性差・層間距離が形へ影響する。
地層を硬い弾性棒と同じものとせず、粘性流動、塑性変形、圧力溶解や層間すべりを考える
\cite{str_fold}。

物理学との接続として、曲げ剛性 \(EI\) の細い弾性棒に軸圧縮力 \(P\) が加わる小変位モデルでは、
\[
 EI\,\frac{d^4y}{dx^4}+P\,\frac{d^2y}{dx^2}=0
\]
となる。\(E\) はヤング率、\(I\) は断面二次モーメント、\(y\) は横変位である。
両端ピンの長さ \(L\) の棒なら最初の固有モードで \(P_{\rm cr}=\pi^2EI/L^2\) となる。
これは層に沿う圧縮が曲げを生むことを示す教材上の類比であり、
自然の粘性地層の臨界荷重や波長をこの式で確定するものではない。

\section{層面すべり・キンク・スランプ}
曲げた複数枚の紙では、外側と内側の弧長差を紙間のすべりが調整する。
地層でも弱い層や境界へ変形が集中する層面すべり褶曲がある。
層間が完全接着した場合と、すべり得る場合は異なるモデルである。
キンクは比較的直線的な翼と局所的な折れ曲がりで姿勢が変わる形式であり、
物性の異方性と関係する\cite{str_fold}。
スランプ褶曲は堆積物の斜面移動に伴う変形である。
移動物の前方停止と後方圧縮は形成の一例であり、すべての過程を一つの停止機構へ固定しない。
令和6年度の強度・剛性率の量的記述とスランプの全称的な説明は追加査読へ残す。

\section{断層の基本：上盤と下盤、相対変位}
\label{theory:structure:faults}
断層は両側の岩体に相対変位を生じた面又は帯である。
傾斜した面の上側が上盤、下側が下盤であり、上盤の相対下降が正断層、上昇が逆断層となる。
横ずれは走向に沿う成分である。向かい側の地塊が右へ動いて見えれば右横ずれ、
左へ動いて見えれば左横ずれという。斜めすべりは複数成分を持つ。
絶対標高や、一つの岩体が東へ動くという記述だけでは運動の分類を確定できない。
条線はすべり軸を記録するが、単一の直線には向きの正負の情報がない。
センスの決定には段差・非対称組織・変位した標識などの追加観察が必要となる。

\section{地塁・地溝と撓曲}
\label{theory:structure:horst}
両側に断層をもつ細長い地塊で、周囲より相対的に高いものが地塁、低いものが地溝である。
同じ正断層群でも、どの地塊を比較するかによって高低が異なる。
原典の「落ち込む地帯を地塁とする」は上下を取り違えた記述である。
基盤の変位が被覆層で広がって現れる場合、地層は線状の切断だけでなく
単斜又は撓曲という折れ曲がりを示す。地下断層と地表の形は常に同じではない。

\section{分岐・低角面・湾曲面}
\label{theory:structure:geometry}
衝上断層で大きく水平移動した岩体がナップ、その下の剥離面がデコルマン又はデタッチメントである。
底の平坦部から傾斜部（ランプ）へ移る上盤の運動は、断層折れ曲がり褶曲をつくり得る。
横ずれ断層帯には、浅部へ枝分かれし深部で主断層へ収束する花束構造が見られる。
これらは移動岩体、すべり面、断層群の断面形態という別々の記載である。
リストリック正断層は、典型的には浅部で急、深部で緩くなる湾曲面である。
「浅部が緩く深部が急」は逆の説明となる。
地すべりの地層境界にも湾曲して深部へ収束する正断層例があり、形態と形成場を分けて使う
\cite{str_listric}。

\section{横ずれの開く場所：プルアパート}
\label{theory:structure:pullapart}
横ずれ断層が一直線でなく、屈曲又は互い違いの区間を持つと、
主断層に沿う運動が局所的な伸張又は圧縮を生む。
開く配置に生じる伸張性の盆地がプルアパート構造である。
反対の配置では圧縮・隆起が生じ得るので、横ずれだからすべて盆地とするのは誤りである。
左右のずれと、段差が開くか閉じるかを別に図示して判定する。

\section{雁行と共役、配列名と変位センス}
\label{theory:structure:conjugate}
雁行は、複数の構造がほぼ平行に斜交しながら階段状に並ぶ幾何学的な配置である。
共役断層は同一の応力場等に対応する異なる方向の群であり、平行な段状配列の別名ではない。
交差する二方向が見えても、形成時期・応力場が違う二群を誤認する可能性がある。
切断関係はイベントの順序の手掛かりとなるが、既存断層の再活動があれば
最後の活動順序と初生形成の順序は一致しない場合がある\cite{str_orientation}。
ミ型・杉型という配列名を、そのまま右横ずれ・左横ずれという運動名へ置き換えない。
熊本地震の大学調査報告は、右横ずれに伴う杉型の地表配列を記録する
\cite{str_en2016}。亀裂の種類、観察方向、図の回転、引きずりの記載を固定して用いる。
原典3問の配列からセンスへの適用条件は追加査読へ残す。

\section{断層岩：形成機構と分類軸}
\label{theory:structure:rocks}
断層岩は、すべり帯で原岩が破砕・変形・融解した記録である。
未固結で細粒基質に富む断層ガウジと、固結した破砕岩のカタクレーサイトを区別する。
後者のプロト・通常・ウルトラ区分には、基質と破砕片の割合を用いる。
一つの分類案では細粒基質が10〜50、50〜90、90〜100\%という区分となる。
分類案や境界の記載方式を確認し、固結度をこの割合の代用にしない。
マイロナイトは結晶塑性等で粒径や配列が変化した岩石であり、
同じ細粒というだけで破砕岩と同一機構にしない\cite{str_faultrocks}。

\section{摩擦溶融と非対称微小構造}
\label{theory:structure:shear}
シュードタキライトは摩擦等の熱で局所的に溶融して急冷した断層岩である。
ガラス質基質、融食された鉱物片等は形成過程の手掛かりであり、
すべての試料で同じ特徴が全部残るとは限らない。
剛体粒子・層状鉱物等に見られる紡錘形のフィッシュや、
斑状粒子の非対称な尾は剪断センスの指標となる。
その判定は線構造に平行で面構造に垂直な断面を観察するなど、方位条件が必要である。
ガウジの粒子回転や粘土による可塑的な挙動は、
鉱物結晶の塑性変形や動的再結晶と区別する\cite{str_shear}。

\section{地下の破壊と地表の表れ}
\label{theory:structure:source}
地震を発生させた地下の破壊面が震源断層、地震時にずれが地表まで達して表れるものが
地表地震断層である。地下の破壊の全範囲が地表に露出するとは限らない。
活断層は最近の地質時代に活動し将来も活動すると考えられる断層である。
地質時代の採用期間には数十万年前以降を重視するものや第四紀中の証拠を含めるものがある。
過去の証拠から将来の可能性を評価することと、特定年月の予知は異なる
\cite{str_source}。
平成30年度の「地震断層」の史料条件は歴史例に限定した呼称の可能性があり、
地表地震断層一般に必要な条件かを査読対象として残す。

\section{活動度：累積した変位を時間で割る}
\label{theory:structure:activity}
長期間の変位量 \(\Delta u\) と期間 \(\Delta t\) から、
\(\bar v=\Delta u/\Delta t\) という平均変位速度を求める。
これは地震時の瞬間すべり速度ではない。
\[
1\,{\rm m}/1000\,{\rm yr}=1\,{\rm mm}/{\rm yr}.
\]
活動度Aは1以上10未満、Bは0.1以上1未満、Cは0.01以上0.1未満m/千年である
\cite{str_activity}。断層運動の成分、累積変位の測定対象、年代誤差を確認する。
この活動度A/B/Cと、一定期間の地震発生確率に付けたランクは別の指標である。
規模と断層長・すべり量の関係は統計的傾向であり、単一の長さから規模を厳密に決めない。

\section{変位地形と、別の過程で作られた地形}
\label{theory:structure:landforms}
閉塞丘は横ずれで移動した尾根等が谷を遮る地形、撓曲崖は層の曲がりとして現れる崖である。
三角末端面は山脚の切断とその後の侵食等を反映し、バルジは圧縮・隆起部の形態となる。
どれも形だけで断層を確定せず、地層対比や累積変位などを合わせる。
モレーンは氷河が運ぶ岩屑の堆積地形である\cite{str_moraine}。
その堆積物が後に断層で変位する場合も、堆積地形としての成因は氷河である。

\section{調査の役割：観察する量を先に決める}
\label{theory:structure:survey}
空中写真は変位地形と広域の連続性、踏査は露頭の地層・岩相・姿勢・断層岩、
トレンチは断層と地層の切断・被覆を観察する。
年代資料を組み合わせて活動履歴を制約する。
磁気探査は岩石の磁化のコントラストを利用して地下の境界を推定するものであり、
磁気異常を断層そのものと同一視しない。
GSJの出題前年の資料は沿岸活断層の連続性・活動性調査に空中磁気探査を含める
\cite{str_coastal2017}。
花粉分析は古植生・古環境や地層の対比を補う。
野坂断層の研究では、花粉・珪藻分析を堆積環境の変遷のために実施している
\cite{str_pollen2014}。直接変位と間接環境資料を分けると役割を理解できる。
平成30年度問23の初回候補5と公式4は、調査目的を省いた相対順位の疑義として保持する。

\section{貫入・浮力と地質構造}
\label{theory:structure:intrusions}
シルは層に沿う板状貫入岩体、ラコリスは屋根を膨らませた貫入岩体である
\cite{str_intrusion}。周囲の地層を曲げることはあっても、ヒンジや背斜などの褶曲用語とは
記載対象が違う。
ダイアピールは低密度等の塑性物質の上昇・貫入による構造である
\cite{str_diapir}。密度差のある流体系では重力による位置エネルギー低下が
上昇の駆動の一つとなる。蓋層の強度、粘性、差応力も形へ影響する。
断層との共存・相互制御を否定せず、「最も関係が薄い」は中心的な形成機構の比較として扱う。

\section{日本列島への接続}
\label{theory:structure:japan}
変成帯は温度・圧力の履歴を、断層は変位の履歴を表す。
領家の高温低圧と三波川の高圧型の対比は岩石学へ接続する\cite{str_ryoke}。
島弧の形成に関与する沈み込みプレートと、島弧を載せるプレートは同じ対象ではない。
方位と「年何cm」という速度には相手又は固定した参照系が必要である。
北海道の弧の衝突、九州の伸張など地域の形成史を、単一の局所断層の分類へ置換しない。

\section{典型的な問われ方と学習の順序}
今回照合した20問には、用語の関連を比べる5問、構造・活動・断層岩等の記述を比べる14問、
姿勢測定の手順1問がある。これは確認した群の内訳であり、全期間の頻度の確定値ではない。
最初に面と線の姿勢、上盤と下盤、褶曲の形態と年代を学ぶ。
次に断層岩の機構と分類軸、活動度と調査目的を結ぶ。
繰り返し問題でも、固結度と割合、地下と地表、平行と直交、地塁と地溝、
配列名と運動名、史料条件の追加という違いを比較する。
全問の分類・解説が揃ってから確認済み統計による頻度分析へ接続する。

\section{過去問}
\subsection{平成25年度 問20}
\label{q:h25-reg-applied-20}

\noindent\textbf{原典：} sources/H25/17\_applied\_science.pdf、PDF物理8頁。

\textbf{公式正答：} 5（single）。

\par\textit{原典本文の校正と独自解説は未作成。分類のみの登録項目。}

\subsection{平成25年度 問21}
\label{q:h25-reg-applied-21}

\noindent\textbf{原典：} sources/H25/17\_applied\_science.pdf、PDF物理8頁。

\textbf{公式正答：} 5（single）。

\par\textit{原典本文の校正と独自解説は未作成。分類のみの登録項目。}

\subsection{平成26年度 問35}
\label{q:h26-reg-applied-35}

\noindent\textbf{原典：} sources/H26/17\_applied\_science.pdf、PDF物理14頁。

\textbf{公式正答：} 2（single）。

\par\textit{原典本文の校正と独自解説は未作成。分類のみの登録項目。}

\subsection{平成30年度 問23}
\label{q:h30-reg-applied-23}

\noindent\textbf{原典：} sources/H30/17\_applied\_science.pdf、PDF物理8頁。

\textbf{公式正答：} 4（single）。

\par\textit{原典本文の校正と独自解説は未作成。分類のみの登録項目。}

\subsection{平成30年度 問26}
\label{q:h30-reg-applied-26}

\noindent\textbf{原典：} sources/H30/17\_applied\_science.pdf、PDF物理9頁。

\textbf{公式正答：} 5（single）。

\par\textit{原典本文の校正と独自解説は未作成。分類のみの登録項目。}

\subsection{令和元年度 問31}
\label{q:r01-reg-applied-31}

\noindent\textbf{原典：} sources/R01/17\_applied\_science.pdf、PDF物理11頁。

\textbf{公式正答：} 2（single）。

\par\textit{原典本文の校正と独自解説は未作成。分類のみの登録項目。}

\subsection{令和元年度再試験 問30}
\label{q:r01-ret-applied-30}

\noindent\textbf{原典：} sources/R01\_retest/17\_applied\_science.pdf、PDF物理10頁。

\textbf{公式正答：} 2（single）。

\par\textit{原典本文の校正と独自解説は未作成。分類のみの登録項目。}

\subsection{令和2年度 問27}
\label{q:r02-reg-applied-27}

\noindent\textbf{原典：} sources/R02/17\_applied\_science.pdf、PDF物理10頁。

\textbf{公式正答：} 4（single）。

\par\textit{原典本文の校正と独自解説は未作成。分類のみの登録項目。}

\subsection{令和2年度 問30}
\label{q:r02-reg-applied-30}

\noindent\textbf{原典：} sources/R02/17\_applied\_science.pdf、PDF物理11頁。

\textbf{公式正答：} 3（single）。

\par\textit{原典本文の校正と独自解説は未作成。分類のみの登録項目。}

\subsection{令和3年度 問27}
\label{q:r03-reg-applied-27}

\noindent\textbf{原典：} sources/R03/17\_applied\_science.pdf、PDF物理10頁。

\textbf{公式正答：} 5（single）。

\par\textit{原典本文の校正と独自解説は未作成。分類のみの登録項目。}

\subsection{令和4年度 問29}
\label{q:r04-reg-applied-29}

\noindent\textbf{原典：} sources/R04/17\_applied\_science.pdf、PDF物理10頁。

\textbf{公式正答：} 5（single）。

\par\textit{原典本文の校正と独自解説は未作成。分類のみの登録項目。}

\subsection{令和4年度 問30}
\label{q:r04-reg-applied-30}

\noindent\textbf{原典：} sources/R04/17\_applied\_science.pdf、PDF物理10頁。

\textbf{公式正答：} 2（single）。

\par\textit{原典本文の校正と独自解説は未作成。分類のみの登録項目。}

\subsection{令和4年度 問31}
\label{q:r04-reg-applied-31}

\noindent\textbf{原典：} sources/R04/17\_applied\_science.pdf、PDF物理11頁。

\textbf{公式正答：} 4（single）。

\par\textit{原典本文の校正と独自解説は未作成。分類のみの登録項目。}

\subsection{令和5年度 問27}
\label{q:r05-reg-applied-27}

\noindent\textbf{原典：} sources/R05/17\_applied\_science.pdf、PDF物理10頁。

\textbf{公式正答：} 4（single）。

\par\textit{原典本文の校正と独自解説は未作成。分類のみの登録項目。}

\subsection{令和5年度 問28}
\label{q:r05-reg-applied-28}

\noindent\textbf{原典：} sources/R05/17\_applied\_science.pdf、PDF物理10頁。

\textbf{公式正答：} 2（single）。

\par\textit{原典本文の校正と独自解説は未作成。分類のみの登録項目。}

\subsection{令和5年度 問29}
\label{q:r05-reg-applied-29}

\noindent\textbf{原典：} sources/R05/17\_applied\_science.pdf、PDF物理10頁。

\textbf{公式正答：} 1（single）。

\par\textit{原典本文の校正と独自解説は未作成。分類のみの登録項目。}

\subsection{令和6年度 問28}
\label{q:r06-reg-applied-28}

\noindent\textbf{原典：} sources/R06/17\_applied\_science.pdf、PDF物理9頁。

\textbf{公式正答：} 5（single）。

\par\textit{原典本文の校正と独自解説は未作成。分類のみの登録項目。}

\subsection{令和6年度 問29}
\label{q:r06-reg-applied-29}

\noindent\textbf{原典：} sources/R06/17\_applied\_science.pdf、PDF物理10頁。

\textbf{公式正答：} 3（single）。

\par\textit{原典本文の校正と独自解説は未作成。分類のみの登録項目。}

\subsection{令和7年度 問29}
\label{q:r07-reg-applied-29}

\noindent\textbf{原典：} sources/R07/17\_applied\_science.pdf、PDF物理11頁。

\textbf{公式正答：} 4（single）。

\par\textit{原典本文の校正と独自解説は未作成。分類のみの登録項目。}

\subsection{令和7年度 問30}
\label{q:r07-reg-applied-30}

\noindent\textbf{原典：} sources/R07/17\_applied\_science.pdf、PDF物理11頁。

\textbf{公式正答：} 5（single）。

\par\textit{原典本文の校正と独自解説は未作成。分類のみの登録項目。}

